\section{Additional evaluation}
\label{sec:additional-evaluation}
\expmark{The tables in this appendix contain experiments still to complete.
They extend the existing results with a common four-benchmark comparison,
independent behavioral evaluation, and controlled mechanism tests.}
The protocols below retain the planned condition-composition extension
described in Appendix~\ref{sec:executable-analysis}.

\subsection{Repair success and behavioral completeness}
\label{sec:heldout-main}
The seven methods share the coding agent, backbone, historical source pool,
target tasks, tools, and total budget. ContextGraph and FAISS Flat retrieve
the same three source IDs; ContextGraph additionally uses their relations to
compose checks. The comparison uses a fixed memory bank. Updating memory
after each task is a separate experiment with the same allowed feedback and
task order for every method.

Independent evaluators specify current requirements on up to 40 untouched
tasks per benchmark. Their checks are sealed before seeing the new patches
and are withheld from the coding agents. Combined success requires official
resolution, the current/historical combinations, and preservation of selected
neighboring behavior. Official success with a failed valid requirement counts
as an omission. Appendix~\ref{sec:heldout-protocol} gives the sampling and
annotation procedure.
\begin{table}[!htbp]
\expcolor
\centering
\caption{Matched first-pass official resolution on four benchmarks.
Cells report resolved/planned and assessed/planned. The last row gives the
paired ContextGraph--Flat difference in percentage points, its 95\% interval,
and the jointly assessed count.}
\label{tab:planned-fourbench}
\small
\begin{tabular}{lcccc}
\toprule
Method & Verified & Related-Lite & LoLBench & DeepSWE \\
Planned tasks & 150 & 99 & 100 & 113 \\
\midrule
No memory & \pending & \pending & \pending & \pending \\
FAISS Flat & \pending & \pending & \pending & \pending \\
ExpeL & \pending & \pending & \pending & \pending \\
ExpeRepair & \pending & \pending & \pending & \pending \\
ReasoningBank & \pending & \pending & \pending & \pending \\
ACE & \pending & \pending & \pending & \pending \\
ContextGraph & \pending & \pending & \pending & \pending \\
\midrule
Graph $-$ Flat (pp, CI) & \pending & \pending & \pending & \pending \\
\bottomrule
\end{tabular}
\end{table}
\begin{table}[!htbp]
\expcolor
\centering
\caption{Combined success on independently specified held-out checks.
Cells report successes/eligible and assessed/eligible. Report sampled and
eligible counts for each benchmark; the final row gives the paired difference
and 95\% interval. Table~\ref{tab:planned-heldout-detail} separates the outcomes.}
\label{tab:planned-heldout}
\small
\begin{tabular}{lcccc}
\toprule
Method & Verified & Related-Lite & LoLBench & DeepSWE \\
\midrule
No memory & \pending & \pending & \pending & \pending \\
FAISS Flat & \pending & \pending & \pending & \pending \\
ExpeL & \pending & \pending & \pending & \pending \\
ExpeRepair & \pending & \pending & \pending & \pending \\
ReasoningBank & \pending & \pending & \pending & \pending \\
ACE & \pending & \pending & \pending & \pending \\
ContextGraph & \pending & \pending & \pending & \pending \\
\midrule
CG $-$ Flat (pp, CI) & \pending & \pending & \pending & \pending \\
\bottomrule
\end{tabular}
\end{table}
\subsection{Separate checks, joint checks, and relation information}
The development study follows the three behavioral questions in
Sections~\ref{sec:content}--\ref{sec:feedback-results}, using 24 tasks and three
repetitions. Content interventions compare how text and executable examples
retain a trigger (Table~\ref{tab:planned-content}). H runs the current
and historical checks separately; I replaces the historical check with a joint
check. Both receive the same history and current program, with two check slots,
the same starting state, and matched execution and feedback budgets.
C versus B tests stored input relations against relations reconstructed
from text, using the same composer. C versus G tests a fixed adapter against
free-form check generation, using the same bindings. F versus C tests the
effect of supplying check results before editing. Table~\ref{tab:planned-graph}
defines each condition; together they test triggers, composition, and feedback.
\begin{table}[!htbp]
\expcolor
\centering
\caption{Development ablations on 24 targets with three repetitions.
B/C share a composer; C/G share bindings. H/I compare separate and joint
checks with the same two slots and feedback schedule. Report mean $\pm$ SD
and paired intervals. Check readiness counts valid delivered checks;
Appendix~\ref{sec:planned-graph} specifies the comparisons.}
\label{tab:planned-graph}
\small
\begin{tabular}{p{0.40\linewidth}cccc}
\toprule
Variant & \shortstack{Check\\ready} & Joint & Combined & Regressions \\
\midrule
A. Flat records; no composition & n/a & \pending & \pending & \pending \\
B. Reconstruct all bindings; shared composer & \pending & \pending & \pending & \pending \\
C. Supplied bindings; same composer & \pending & \pending & \pending & \pending \\
D. C with input bindings reconstructed & \pending & \pending & \pending & \pending \\
E. C with operation bindings reconstructed & \pending & \pending & \pending & \pending \\
F. C + supplied execution observations & \pending & \pending & \pending & \pending \\
G. C's bindings; free-form composer & \pending & \pending & \pending & \pending \\
\midrule
H. Current + separate historical check & \pending & \pending & \pending & \pending \\
I. Current + joint check; same feedback & \pending & \pending & \pending & \pending \\
\bottomrule
\end{tabular}
\end{table}
\subsection{Where the method helps and what it costs}
Before running the methods, group tasks by whether the adapters support their
operations: supported, outside scope, or uncertain. Compare methods on the
same tasks within each group, including failed preparation and text fallback.
This distinguishes a lack of coverage from a failed repair.
Cost includes source construction, program preparation, retrieval, applicability
assessment, check construction and execution, repair attempts, and verification.
The retry and amortization tables report how these costs change as the bank
is reused.
\begin{table}[!htbp]
\expcolor
\centering
\caption{Coverage and cost by benchmark. Preparation, adapter support,
and joint delivery use all planned targets. Combined differences use the
independent behavioral sample, grouped by support before any run.
USD per solution includes preparation and all repair attempts.}
\label{tab:planned-main-coverage}
\small
\begin{tabular}{lccccc}
\toprule
Benchmark & $P_q$ ready & Supported & Joint delivery & \shortstack{CG--Flat\\Combined (pp)} & \shortstack{CG/Flat\\USD/solve} \\
\midrule
Verified150 & \pending & \pending & \pending & \pending & \pending \\
Related-Lite99 & \pending & \pending & \pending & \pending & \pending \\
LoLBench100 & \pending & \pending & \pending & \pending & \pending \\
DeepSWE113 & \pending & \pending & \pending & \pending & \pending \\
\bottomrule
\end{tabular}
\end{table}
